\documentclass[12pt,a4paper]{article}
\usepackage[top=1.8cm,bottom=1.8cm,left=1.8cm,right=1.8cm]{geometry}
\usepackage{fontspec}
\setmainfont{TeX Gyre Termes}
\usepackage{amsmath,amssymb}
\usepackage{fancyhdr}
\usepackage{tasks}
\usepackage{tikz}
\settasks{label=\textbf{\Alph*.},label-width=1.5em,item-indent=2em}
\setlength{\parindent}{0pt}
\setlength{\parskip}{3pt}
\pagestyle{fancy}\fancyhf{}\renewcommand{\headrulewidth}{0pt}
\fancyhead[L]{\small Ngân hàng bài tập Toán 8 -- 8H2-C6-B24-D02}
\cfoot{\small Trang \thepage}
\newif\ifloigiai
\loigiaitrue
\newenvironment{cauhoi}[2]{%
  \par\medskip\noindent\rule{\textwidth}{0.4pt}\par\nopagebreak
  \noindent\textbf{[#1]}\hfill{\scriptsize\ttfamily #2}\par\nopagebreak\smallskip}{\par}
\newcommand{\loigiai}[1]{\ifloigiai\par\smallskip\noindent\textit{Lời giải.} #1\par\fi}
\newcommand{\dapan}[1]{\ifloigiai\par\smallskip\noindent\textit{Đáp án:} \textbf{#1}\par\fi}
\begin{document}
\begin{center}
{\Large\textbf{NGÂN HÀNG BÀI TẬP TOÁN 8}}\\[3pt]
\textbf{Chương 6. Phân thức đại số}\\[2pt]
\textbf{Bài 24. Phép nhân và phép chia phân thức đại số}\\[2pt]
\textbf{Dạng D02: Biểu thức tổng hợp, tìm x nguyên để giá trị nguyên}\\[2pt]
\textit{Mã dạng: 8H2-C6-B24-D02 \quad 50 bài tự luận + 50 câu trắc nghiệm \quad (VD: 25, VDC: 25)}
\end{center}
\textbf{Quy ước mã câu:} \textit{khối, học kì, chương, bài, dạng, loại (TL/TN), mức độ, số thứ tự}.\par Ví dụ: \texttt{8H2-C6-B24-D02-TL-NB-001}.\par Để in bản cho học sinh, đặt \verb|\loigiaifalse| ngay sau \verb|\loigiaitrue|.
\section*{Phần I. Tự luận}
\begin{cauhoi}{Bài 1}{8H2-C6-B24-D02-TL-VD-001}
Cho $P=\dfrac{x+3}{x+4}$. Tìm các số nguyên $x$ để $P$ có giá trị nguyên.
\loigiai{Viết $P=\dfrac{(x+4)-1}{x+4}=1+\dfrac{-1}{x+4}$. $P$ nguyên khi $x+4$ là ước của $-1$, tức $x+4\in\{-1; 1\}$. Suy ra $x\in\{-5; -3\}$.}
\end{cauhoi}
\begin{cauhoi}{Bài 2}{8H2-C6-B24-D02-TL-VD-002}
Cho $P=\dfrac{x+6}{x+7}$. Tìm các số nguyên $x$ để $P$ có giá trị nguyên.
\loigiai{Viết $P=\dfrac{(x+7)-1}{x+7}=1+\dfrac{-1}{x+7}$. $P$ nguyên khi $x+7$ là ước của $-1$, tức $x+7\in\{-1; 1\}$. Suy ra $x\in\{-8; -6\}$.}
\end{cauhoi}
\begin{cauhoi}{Bài 3}{8H2-C6-B24-D02-TL-VD-003}
Cho $P=\dfrac{x+4}{x+1}$. Tìm các số nguyên $x$ để $P$ có giá trị nguyên.
\loigiai{Viết $P=\dfrac{(x+1)+3}{x+1}=1+\dfrac{3}{x+1}$. $P$ nguyên khi $x+1$ là ước của $3$, tức $x+1\in\{-3; -1; 1; 3\}$. Suy ra $x\in\{-4; -2; 0; 2\}$.}
\end{cauhoi}
\begin{cauhoi}{Bài 4}{8H2-C6-B24-D02-TL-VD-004}
Cho $P=\dfrac{x+3}{x+1}$. Tìm các số nguyên $x$ để $P$ có giá trị nguyên.
\loigiai{Viết $P=\dfrac{(x+1)+2}{x+1}=1+\dfrac{2}{x+1}$. $P$ nguyên khi $x+1$ là ước của $2$, tức $x+1\in\{-2; -1; 1; 2\}$. Suy ra $x\in\{-3; -2; 0; 1\}$.}
\end{cauhoi}
\begin{cauhoi}{Bài 5}{8H2-C6-B24-D02-TL-VD-005}
Cho $P=\dfrac{x+5}{x+3}$. Tìm các số nguyên $x$ để $P$ có giá trị nguyên.
\loigiai{Viết $P=\dfrac{(x+3)+2}{x+3}=1+\dfrac{2}{x+3}$. $P$ nguyên khi $x+3$ là ước của $2$, tức $x+3\in\{-2; -1; 1; 2\}$. Suy ra $x\in\{-5; -4; -2; -1\}$.}
\end{cauhoi}
\begin{cauhoi}{Bài 6}{8H2-C6-B24-D02-TL-VD-006}
Cho $P=\dfrac{x+5}{x+4}$. Tìm các số nguyên $x$ để $P$ có giá trị nguyên.
\loigiai{Viết $P=\dfrac{(x+4)+1}{x+4}=1+\dfrac{1}{x+4}$. $P$ nguyên khi $x+4$ là ước của $1$, tức $x+4\in\{-1; 1\}$. Suy ra $x\in\{-5; -3\}$.}
\end{cauhoi}
\begin{cauhoi}{Bài 7}{8H2-C6-B24-D02-TL-VD-007}
Cho $P=\dfrac{x+5}{x+1}$. Tìm các số nguyên $x$ để $P$ có giá trị nguyên.
\loigiai{Viết $P=\dfrac{(x+1)+4}{x+1}=1+\dfrac{4}{x+1}$. $P$ nguyên khi $x+1$ là ước của $4$, tức $x+1\in\{-4; -2; -1; 1; 2; 4\}$. Suy ra $x\in\{-5; -3; -2; 0; 1; 3\}$.}
\end{cauhoi}
\begin{cauhoi}{Bài 8}{8H2-C6-B24-D02-TL-VD-008}
Cho $P=\dfrac{x+2}{x+8}$. Tìm các số nguyên $x$ để $P$ có giá trị nguyên.
\loigiai{Viết $P=\dfrac{(x+8)-6}{x+8}=1+\dfrac{-6}{x+8}$. $P$ nguyên khi $x+8$ là ước của $-6$, tức $x+8\in\{-6; -3; -2; -1; 1; 2; 3; 6\}$. Suy ra $x\in\{-14; -11; -10; -9; -7; -6; -5; -2\}$.}
\end{cauhoi}
\begin{cauhoi}{Bài 9}{8H2-C6-B24-D02-TL-VD-009}
Cho $P=\dfrac{x+4}{x+7}$. Tìm các số nguyên $x$ để $P$ có giá trị nguyên.
\loigiai{Viết $P=\dfrac{(x+7)-3}{x+7}=1+\dfrac{-3}{x+7}$. $P$ nguyên khi $x+7$ là ước của $-3$, tức $x+7\in\{-3; -1; 1; 3\}$. Suy ra $x\in\{-10; -8; -6; -4\}$.}
\end{cauhoi}
\begin{cauhoi}{Bài 10}{8H2-C6-B24-D02-TL-VD-010}
Cho $P=\dfrac{x+2}{x+3}$. Tìm các số nguyên $x$ để $P$ có giá trị nguyên.
\loigiai{Viết $P=\dfrac{(x+3)-1}{x+3}=1+\dfrac{-1}{x+3}$. $P$ nguyên khi $x+3$ là ước của $-1$, tức $x+3\in\{-1; 1\}$. Suy ra $x\in\{-4; -2\}$.}
\end{cauhoi}
\begin{cauhoi}{Bài 11}{8H2-C6-B24-D02-TL-VD-011}
Cho $P=\dfrac{x+1}{x+3}$. Tìm các số nguyên $x$ để $P$ có giá trị nguyên.
\loigiai{Viết $P=\dfrac{(x+3)-2}{x+3}=1+\dfrac{-2}{x+3}$. $P$ nguyên khi $x+3$ là ước của $-2$, tức $x+3\in\{-2; -1; 1; 2\}$. Suy ra $x\in\{-5; -4; -2; -1\}$.}
\end{cauhoi}
\begin{cauhoi}{Bài 12}{8H2-C6-B24-D02-TL-VD-012}
Cho $P=\dfrac{x+8}{x+7}$. Tìm các số nguyên $x$ để $P$ có giá trị nguyên.
\loigiai{Viết $P=\dfrac{(x+7)+1}{x+7}=1+\dfrac{1}{x+7}$. $P$ nguyên khi $x+7$ là ước của $1$, tức $x+7\in\{-1; 1\}$. Suy ra $x\in\{-8; -6\}$.}
\end{cauhoi}
\begin{cauhoi}{Bài 13}{8H2-C6-B24-D02-TL-VD-013}
Cho $P=\dfrac{x+8}{x+2}$. Tìm các số nguyên $x$ để $P$ có giá trị nguyên.
\loigiai{Viết $P=\dfrac{(x+2)+6}{x+2}=1+\dfrac{6}{x+2}$. $P$ nguyên khi $x+2$ là ước của $6$, tức $x+2\in\{-6; -3; -2; -1; 1; 2; 3; 6\}$. Suy ra $x\in\{-8; -5; -4; -3; -1; 0; 1; 4\}$.}
\end{cauhoi}
\begin{cauhoi}{Bài 14}{8H2-C6-B24-D02-TL-VD-014}
Cho $P=\dfrac{x+4}{x+6}$. Tìm các số nguyên $x$ để $P$ có giá trị nguyên.
\loigiai{Viết $P=\dfrac{(x+6)-2}{x+6}=1+\dfrac{-2}{x+6}$. $P$ nguyên khi $x+6$ là ước của $-2$, tức $x+6\in\{-2; -1; 1; 2\}$. Suy ra $x\in\{-8; -7; -5; -4\}$.}
\end{cauhoi}
\begin{cauhoi}{Bài 15}{8H2-C6-B24-D02-TL-VD-015}
Cho $P=\dfrac{x+4}{x+5}$. Tìm các số nguyên $x$ để $P$ có giá trị nguyên.
\loigiai{Viết $P=\dfrac{(x+5)-1}{x+5}=1+\dfrac{-1}{x+5}$. $P$ nguyên khi $x+5$ là ước của $-1$, tức $x+5\in\{-1; 1\}$. Suy ra $x\in\{-6; -4\}$.}
\end{cauhoi}
\begin{cauhoi}{Bài 16}{8H2-C6-B24-D02-TL-VD-016}
Cho $P=\dfrac{x+3}{x+6}$. Tìm các số nguyên $x$ để $P$ có giá trị nguyên.
\loigiai{Viết $P=\dfrac{(x+6)-3}{x+6}=1+\dfrac{-3}{x+6}$. $P$ nguyên khi $x+6$ là ước của $-3$, tức $x+6\in\{-3; -1; 1; 3\}$. Suy ra $x\in\{-9; -7; -5; -3\}$.}
\end{cauhoi}
\begin{cauhoi}{Bài 17}{8H2-C6-B24-D02-TL-VD-017}
Cho $P=\dfrac{x+8}{x+3}$. Tìm các số nguyên $x$ để $P$ có giá trị nguyên.
\loigiai{Viết $P=\dfrac{(x+3)+5}{x+3}=1+\dfrac{5}{x+3}$. $P$ nguyên khi $x+3$ là ước của $5$, tức $x+3\in\{-5; -1; 1; 5\}$. Suy ra $x\in\{-8; -4; -2; 2\}$.}
\end{cauhoi}
\begin{cauhoi}{Bài 18}{8H2-C6-B24-D02-TL-VD-018}
Cho $P=\dfrac{x+1}{x+2}$. Tìm các số nguyên $x$ để $P$ có giá trị nguyên.
\loigiai{Viết $P=\dfrac{(x+2)-1}{x+2}=1+\dfrac{-1}{x+2}$. $P$ nguyên khi $x+2$ là ước của $-1$, tức $x+2\in\{-1; 1\}$. Suy ra $x\in\{-3; -1\}$.}
\end{cauhoi}
\begin{cauhoi}{Bài 19}{8H2-C6-B24-D02-TL-VD-019}
Cho $P=\dfrac{x+7}{x+2}$. Tìm các số nguyên $x$ để $P$ có giá trị nguyên.
\loigiai{Viết $P=\dfrac{(x+2)+5}{x+2}=1+\dfrac{5}{x+2}$. $P$ nguyên khi $x+2$ là ước của $5$, tức $x+2\in\{-5; -1; 1; 5\}$. Suy ra $x\in\{-7; -3; -1; 3\}$.}
\end{cauhoi}
\begin{cauhoi}{Bài 20}{8H2-C6-B24-D02-TL-VD-020}
Cho $P=\dfrac{x+6}{x+8}$. Tìm các số nguyên $x$ để $P$ có giá trị nguyên.
\loigiai{Viết $P=\dfrac{(x+8)-2}{x+8}=1+\dfrac{-2}{x+8}$. $P$ nguyên khi $x+8$ là ước của $-2$, tức $x+8\in\{-2; -1; 1; 2\}$. Suy ra $x\in\{-10; -9; -7; -6\}$.}
\end{cauhoi}
\begin{cauhoi}{Bài 21}{8H2-C6-B24-D02-TL-VD-021}
Cho $P=\dfrac{x+3}{x+7}$. Tìm các số nguyên $x$ để $P$ có giá trị nguyên.
\loigiai{Viết $P=\dfrac{(x+7)-4}{x+7}=1+\dfrac{-4}{x+7}$. $P$ nguyên khi $x+7$ là ước của $-4$, tức $x+7\in\{-4; -2; -1; 1; 2; 4\}$. Suy ra $x\in\{-11; -9; -8; -6; -5; -3\}$.}
\end{cauhoi}
\begin{cauhoi}{Bài 22}{8H2-C6-B24-D02-TL-VD-022}
Cho $P=\dfrac{x+7}{x+4}$. Tìm các số nguyên $x$ để $P$ có giá trị nguyên.
\loigiai{Viết $P=\dfrac{(x+4)+3}{x+4}=1+\dfrac{3}{x+4}$. $P$ nguyên khi $x+4$ là ước của $3$, tức $x+4\in\{-3; -1; 1; 3\}$. Suy ra $x\in\{-7; -5; -3; -1\}$.}
\end{cauhoi}
\begin{cauhoi}{Bài 23}{8H2-C6-B24-D02-TL-VD-023}
Cho $P=\dfrac{x+7}{x+8}$. Tìm các số nguyên $x$ để $P$ có giá trị nguyên.
\loigiai{Viết $P=\dfrac{(x+8)-1}{x+8}=1+\dfrac{-1}{x+8}$. $P$ nguyên khi $x+8$ là ước của $-1$, tức $x+8\in\{-1; 1\}$. Suy ra $x\in\{-9; -7\}$.}
\end{cauhoi}
\begin{cauhoi}{Bài 24}{8H2-C6-B24-D02-TL-VD-024}
Cho $P=\dfrac{x+5}{x+7}$. Tìm các số nguyên $x$ để $P$ có giá trị nguyên.
\loigiai{Viết $P=\dfrac{(x+7)-2}{x+7}=1+\dfrac{-2}{x+7}$. $P$ nguyên khi $x+7$ là ước của $-2$, tức $x+7\in\{-2; -1; 1; 2\}$. Suy ra $x\in\{-9; -8; -6; -5\}$.}
\end{cauhoi}
\begin{cauhoi}{Bài 25}{8H2-C6-B24-D02-TL-VD-025}
Cho $P=\dfrac{x+4}{x+2}$. Tìm các số nguyên $x$ để $P$ có giá trị nguyên.
\loigiai{Viết $P=\dfrac{(x+2)+2}{x+2}=1+\dfrac{2}{x+2}$. $P$ nguyên khi $x+2$ là ước của $2$, tức $x+2\in\{-2; -1; 1; 2\}$. Suy ra $x\in\{-4; -3; -1; 0\}$.}
\end{cauhoi}
\begin{cauhoi}{Bài 26}{8H2-C6-B24-D02-TL-VDC-026}
Cho $P=\dfrac{x+1}{x+4}$. Tìm các số nguyên $x$ để $P$ có giá trị nguyên.
\loigiai{Viết $P=\dfrac{(x+4)-3}{x+4}=1+\dfrac{-3}{x+4}$. $P$ nguyên khi $x+4$ là ước của $-3$, tức $x+4\in\{-3; -1; 1; 3\}$. Suy ra $x\in\{-7; -5; -3; -1\}$.}
\end{cauhoi}
\begin{cauhoi}{Bài 27}{8H2-C6-B24-D02-TL-VDC-027}
Cho $P=\dfrac{x+8}{x+6}$. Tìm các số nguyên $x$ để $P$ có giá trị nguyên.
\loigiai{Viết $P=\dfrac{(x+6)+2}{x+6}=1+\dfrac{2}{x+6}$. $P$ nguyên khi $x+6$ là ước của $2$, tức $x+6\in\{-2; -1; 1; 2\}$. Suy ra $x\in\{-8; -7; -5; -4\}$.}
\end{cauhoi}
\begin{cauhoi}{Bài 28}{8H2-C6-B24-D02-TL-VDC-028}
Cho $P=\dfrac{x+1}{x+8}$. Tìm các số nguyên $x$ để $P$ có giá trị nguyên.
\loigiai{Viết $P=\dfrac{(x+8)-7}{x+8}=1+\dfrac{-7}{x+8}$. $P$ nguyên khi $x+8$ là ước của $-7$, tức $x+8\in\{-7; -1; 1; 7\}$. Suy ra $x\in\{-15; -9; -7; -1\}$.}
\end{cauhoi}
\begin{cauhoi}{Bài 29}{8H2-C6-B24-D02-TL-VDC-029}
Cho $P=\dfrac{x+5}{x+2}$. Tìm các số nguyên $x$ để $P$ có giá trị nguyên.
\loigiai{Viết $P=\dfrac{(x+2)+3}{x+2}=1+\dfrac{3}{x+2}$. $P$ nguyên khi $x+2$ là ước của $3$, tức $x+2\in\{-3; -1; 1; 3\}$. Suy ra $x\in\{-5; -3; -1; 1\}$.}
\end{cauhoi}
\begin{cauhoi}{Bài 30}{8H2-C6-B24-D02-TL-VDC-030}
Cho $P=\dfrac{x+1}{x+7}$. Tìm các số nguyên $x$ để $P$ có giá trị nguyên.
\loigiai{Viết $P=\dfrac{(x+7)-6}{x+7}=1+\dfrac{-6}{x+7}$. $P$ nguyên khi $x+7$ là ước của $-6$, tức $x+7\in\{-6; -3; -2; -1; 1; 2; 3; 6\}$. Suy ra $x\in\{-13; -10; -9; -8; -6; -5; -4; -1\}$.}
\end{cauhoi}
\begin{cauhoi}{Bài 31}{8H2-C6-B24-D02-TL-VDC-031}
Cho $P=\dfrac{x+7}{x+6}$. Tìm các số nguyên $x$ để $P$ có giá trị nguyên.
\loigiai{Viết $P=\dfrac{(x+6)+1}{x+6}=1+\dfrac{1}{x+6}$. $P$ nguyên khi $x+6$ là ước của $1$, tức $x+6\in\{-1; 1\}$. Suy ra $x\in\{-7; -5\}$.}
\end{cauhoi}
\begin{cauhoi}{Bài 32}{8H2-C6-B24-D02-TL-VDC-032}
Cho $P=\dfrac{x+5}{x+6}$. Tìm các số nguyên $x$ để $P$ có giá trị nguyên.
\loigiai{Viết $P=\dfrac{(x+6)-1}{x+6}=1+\dfrac{-1}{x+6}$. $P$ nguyên khi $x+6$ là ước của $-1$, tức $x+6\in\{-1; 1\}$. Suy ra $x\in\{-7; -5\}$.}
\end{cauhoi}
\begin{cauhoi}{Bài 33}{8H2-C6-B24-D02-TL-VDC-033}
Cho $P=\dfrac{x+4}{x+3}$. Tìm các số nguyên $x$ để $P$ có giá trị nguyên.
\loigiai{Viết $P=\dfrac{(x+3)+1}{x+3}=1+\dfrac{1}{x+3}$. $P$ nguyên khi $x+3$ là ước của $1$, tức $x+3\in\{-1; 1\}$. Suy ra $x\in\{-4; -2\}$.}
\end{cauhoi}
\begin{cauhoi}{Bài 34}{8H2-C6-B24-D02-TL-VDC-034}
Cho $P=\dfrac{x+2}{x+4}$. Tìm các số nguyên $x$ để $P$ có giá trị nguyên.
\loigiai{Viết $P=\dfrac{(x+4)-2}{x+4}=1+\dfrac{-2}{x+4}$. $P$ nguyên khi $x+4$ là ước của $-2$, tức $x+4\in\{-2; -1; 1; 2\}$. Suy ra $x\in\{-6; -5; -3; -2\}$.}
\end{cauhoi}
\begin{cauhoi}{Bài 35}{8H2-C6-B24-D02-TL-VDC-035}
Cho $P=\dfrac{x+8}{x+4}$. Tìm các số nguyên $x$ để $P$ có giá trị nguyên.
\loigiai{Viết $P=\dfrac{(x+4)+4}{x+4}=1+\dfrac{4}{x+4}$. $P$ nguyên khi $x+4$ là ước của $4$, tức $x+4\in\{-4; -2; -1; 1; 2; 4\}$. Suy ra $x\in\{-8; -6; -5; -3; -2; 0\}$.}
\end{cauhoi}
\begin{cauhoi}{Bài 36}{8H2-C6-B24-D02-TL-VDC-036}
Cho $P=\dfrac{x+6}{x+1}$. Tìm các số nguyên $x$ để $P$ có giá trị nguyên.
\loigiai{Viết $P=\dfrac{(x+1)+5}{x+1}=1+\dfrac{5}{x+1}$. $P$ nguyên khi $x+1$ là ước của $5$, tức $x+1\in\{-5; -1; 1; 5\}$. Suy ra $x\in\{-6; -2; 0; 4\}$.}
\end{cauhoi}
\begin{cauhoi}{Bài 37}{8H2-C6-B24-D02-TL-VDC-037}
Cho $P=\dfrac{x+1}{x+5}$. Tìm các số nguyên $x$ để $P$ có giá trị nguyên.
\loigiai{Viết $P=\dfrac{(x+5)-4}{x+5}=1+\dfrac{-4}{x+5}$. $P$ nguyên khi $x+5$ là ước của $-4$, tức $x+5\in\{-4; -2; -1; 1; 2; 4\}$. Suy ra $x\in\{-9; -7; -6; -4; -3; -1\}$.}
\end{cauhoi}
\begin{cauhoi}{Bài 38}{8H2-C6-B24-D02-TL-VDC-038}
Cho $P=\dfrac{x+3}{x+5}$. Tìm các số nguyên $x$ để $P$ có giá trị nguyên.
\loigiai{Viết $P=\dfrac{(x+5)-2}{x+5}=1+\dfrac{-2}{x+5}$. $P$ nguyên khi $x+5$ là ước của $-2$, tức $x+5\in\{-2; -1; 1; 2\}$. Suy ra $x\in\{-7; -6; -4; -3\}$.}
\end{cauhoi}
\begin{cauhoi}{Bài 39}{8H2-C6-B24-D02-TL-VDC-039}
Cho $P=\dfrac{x+8}{x+5}$. Tìm các số nguyên $x$ để $P$ có giá trị nguyên.
\loigiai{Viết $P=\dfrac{(x+5)+3}{x+5}=1+\dfrac{3}{x+5}$. $P$ nguyên khi $x+5$ là ước của $3$, tức $x+5\in\{-3; -1; 1; 3\}$. Suy ra $x\in\{-8; -6; -4; -2\}$.}
\end{cauhoi}
\begin{cauhoi}{Bài 40}{8H2-C6-B24-D02-TL-VDC-040}
Cho $P=\dfrac{x+4}{x+8}$. Tìm các số nguyên $x$ để $P$ có giá trị nguyên.
\loigiai{Viết $P=\dfrac{(x+8)-4}{x+8}=1+\dfrac{-4}{x+8}$. $P$ nguyên khi $x+8$ là ước của $-4$, tức $x+8\in\{-4; -2; -1; 1; 2; 4\}$. Suy ra $x\in\{-12; -10; -9; -7; -6; -4\}$.}
\end{cauhoi}
\begin{cauhoi}{Bài 41}{8H2-C6-B24-D02-TL-VDC-041}
Cho $P=\dfrac{x+7}{x+5}$. Tìm các số nguyên $x$ để $P$ có giá trị nguyên.
\loigiai{Viết $P=\dfrac{(x+5)+2}{x+5}=1+\dfrac{2}{x+5}$. $P$ nguyên khi $x+5$ là ước của $2$, tức $x+5\in\{-2; -1; 1; 2\}$. Suy ra $x\in\{-7; -6; -4; -3\}$.}
\end{cauhoi}
\begin{cauhoi}{Bài 42}{8H2-C6-B24-D02-TL-VDC-042}
Cho $P=\dfrac{x+6}{x+3}$. Tìm các số nguyên $x$ để $P$ có giá trị nguyên.
\loigiai{Viết $P=\dfrac{(x+3)+3}{x+3}=1+\dfrac{3}{x+3}$. $P$ nguyên khi $x+3$ là ước của $3$, tức $x+3\in\{-3; -1; 1; 3\}$. Suy ra $x\in\{-6; -4; -2; 0\}$.}
\end{cauhoi}
\begin{cauhoi}{Bài 43}{8H2-C6-B24-D02-TL-VDC-043}
Cho $P=\dfrac{x+2}{x+5}$. Tìm các số nguyên $x$ để $P$ có giá trị nguyên.
\loigiai{Viết $P=\dfrac{(x+5)-3}{x+5}=1+\dfrac{-3}{x+5}$. $P$ nguyên khi $x+5$ là ước của $-3$, tức $x+5\in\{-3; -1; 1; 3\}$. Suy ra $x\in\{-8; -6; -4; -2\}$.}
\end{cauhoi}
\begin{cauhoi}{Bài 44}{8H2-C6-B24-D02-TL-VDC-044}
Cho $P=\dfrac{x+2}{x+7}$. Tìm các số nguyên $x$ để $P$ có giá trị nguyên.
\loigiai{Viết $P=\dfrac{(x+7)-5}{x+7}=1+\dfrac{-5}{x+7}$. $P$ nguyên khi $x+7$ là ước của $-5$, tức $x+7\in\{-5; -1; 1; 5\}$. Suy ra $x\in\{-12; -8; -6; -2\}$.}
\end{cauhoi}
\begin{cauhoi}{Bài 45}{8H2-C6-B24-D02-TL-VDC-045}
Cho $P=\dfrac{x+7}{x+1}$. Tìm các số nguyên $x$ để $P$ có giá trị nguyên.
\loigiai{Viết $P=\dfrac{(x+1)+6}{x+1}=1+\dfrac{6}{x+1}$. $P$ nguyên khi $x+1$ là ước của $6$, tức $x+1\in\{-6; -3; -2; -1; 1; 2; 3; 6\}$. Suy ra $x\in\{-7; -4; -3; -2; 0; 1; 2; 5\}$.}
\end{cauhoi}
\begin{cauhoi}{Bài 46}{8H2-C6-B24-D02-TL-VDC-046}
Cho $P=\dfrac{x+1}{x+6}$. Tìm các số nguyên $x$ để $P$ có giá trị nguyên.
\loigiai{Viết $P=\dfrac{(x+6)-5}{x+6}=1+\dfrac{-5}{x+6}$. $P$ nguyên khi $x+6$ là ước của $-5$, tức $x+6\in\{-5; -1; 1; 5\}$. Suy ra $x\in\{-11; -7; -5; -1\}$.}
\end{cauhoi}
\begin{cauhoi}{Bài 47}{8H2-C6-B24-D02-TL-VDC-047}
Cho $P=\dfrac{x+8}{x+1}$. Tìm các số nguyên $x$ để $P$ có giá trị nguyên.
\loigiai{Viết $P=\dfrac{(x+1)+7}{x+1}=1+\dfrac{7}{x+1}$. $P$ nguyên khi $x+1$ là ước của $7$, tức $x+1\in\{-7; -1; 1; 7\}$. Suy ra $x\in\{-8; -2; 0; 6\}$.}
\end{cauhoi}
\begin{cauhoi}{Bài 48}{8H2-C6-B24-D02-TL-VDC-048}
Cho $P=\dfrac{x+6}{x+4}$. Tìm các số nguyên $x$ để $P$ có giá trị nguyên.
\loigiai{Viết $P=\dfrac{(x+4)+2}{x+4}=1+\dfrac{2}{x+4}$. $P$ nguyên khi $x+4$ là ước của $2$, tức $x+4\in\{-2; -1; 1; 2\}$. Suy ra $x\in\{-6; -5; -3; -2\}$.}
\end{cauhoi}
\begin{cauhoi}{Bài 49}{8H2-C6-B24-D02-TL-VDC-049}
Cho $P=\dfrac{x+5}{x+8}$. Tìm các số nguyên $x$ để $P$ có giá trị nguyên.
\loigiai{Viết $P=\dfrac{(x+8)-3}{x+8}=1+\dfrac{-3}{x+8}$. $P$ nguyên khi $x+8$ là ước của $-3$, tức $x+8\in\{-3; -1; 1; 3\}$. Suy ra $x\in\{-11; -9; -7; -5\}$.}
\end{cauhoi}
\begin{cauhoi}{Bài 50}{8H2-C6-B24-D02-TL-VDC-050}
Cho $P=\dfrac{x+3}{x+2}$. Tìm các số nguyên $x$ để $P$ có giá trị nguyên.
\loigiai{Viết $P=\dfrac{(x+2)+1}{x+2}=1+\dfrac{1}{x+2}$. $P$ nguyên khi $x+2$ là ước của $1$, tức $x+2\in\{-1; 1\}$. Suy ra $x\in\{-3; -1\}$.}
\end{cauhoi}
\clearpage\section*{Phần II. Trắc nghiệm}
\begin{cauhoi}{Câu 1}{8H2-C6-B24-D02-TN-VD-001}
Cho $P=\dfrac{x+3}{x+4}$. Tìm các số nguyên $x$ để $P$ có giá trị nguyên.
\begin{tasks}(2)
\task $\{-1; 1\}$
\task $\{-5\}$
\task $\{-3\}$
\task $\{-5; -3\}$
\end{tasks}
\dapan{D}
\end{cauhoi}
\begin{cauhoi}{Câu 2}{8H2-C6-B24-D02-TN-VD-002}
Cho $P=\dfrac{x+6}{x+7}$. Tìm các số nguyên $x$ để $P$ có giá trị nguyên.
\begin{tasks}(2)
\task $\{-8; -6\}$
\task $\{-6\}$
\task $\{-8\}$
\task $\{-1; 1\}$
\end{tasks}
\dapan{A}
\end{cauhoi}
\begin{cauhoi}{Câu 3}{8H2-C6-B24-D02-TN-VD-003}
Cho $P=\dfrac{x+4}{x+1}$. Tìm các số nguyên $x$ để $P$ có giá trị nguyên.
\begin{tasks}(2)
\task $\{-4; -2; 0\}$
\task $\{-4; -2; 0; 2\}$
\task $\{-2; 0; 2\}$
\task $\{-3; -1; 1; 3\}$
\end{tasks}
\dapan{B}
\end{cauhoi}
\begin{cauhoi}{Câu 4}{8H2-C6-B24-D02-TN-VD-004}
Cho $P=\dfrac{x+3}{x+1}$. Tìm các số nguyên $x$ để $P$ có giá trị nguyên.
\begin{tasks}(2)
\task $\{-2; 0; 1\}$
\task $\{-3; -2; 0; 1\}$
\task $\{-3; -2; 0\}$
\task $\{-2; -1; 1; 2\}$
\end{tasks}
\dapan{B}
\end{cauhoi}
\begin{cauhoi}{Câu 5}{8H2-C6-B24-D02-TN-VD-005}
Cho $P=\dfrac{x+5}{x+3}$. Tìm các số nguyên $x$ để $P$ có giá trị nguyên.
\begin{tasks}(2)
\task $\{-4; -2; -1\}$
\task $\{-2; -1; 1; 2\}$
\task $\{-5; -4; -2; -1\}$
\task $\{-5; -4; -2\}$
\end{tasks}
\dapan{C}
\end{cauhoi}
\begin{cauhoi}{Câu 6}{8H2-C6-B24-D02-TN-VD-006}
Cho $P=\dfrac{x+5}{x+4}$. Tìm các số nguyên $x$ để $P$ có giá trị nguyên.
\begin{tasks}(2)
\task $\{-3\}$
\task $\{-5\}$
\task $\{-1; 1\}$
\task $\{-5; -3\}$
\end{tasks}
\dapan{D}
\end{cauhoi}
\begin{cauhoi}{Câu 7}{8H2-C6-B24-D02-TN-VD-007}
Cho $P=\dfrac{x+5}{x+1}$. Tìm các số nguyên $x$ để $P$ có giá trị nguyên.
\begin{tasks}(2)
\task $\{-4; -2; -1; 1; 2; 4\}$
\task $\{-5; -3; -2; 0; 1\}$
\task $\{-3; -2; 0; 1; 3\}$
\task $\{-5; -3; -2; 0; 1; 3\}$
\end{tasks}
\dapan{D}
\end{cauhoi}
\begin{cauhoi}{Câu 8}{8H2-C6-B24-D02-TN-VD-008}
Cho $P=\dfrac{x+2}{x+8}$. Tìm các số nguyên $x$ để $P$ có giá trị nguyên.
\begin{tasks}(2)
\task $\{-14; -11; -10; -9; -7; -6; -5; -2\}$
\task $\{-11; -10; -9; -7; -6; -5; -2\}$
\task $\{-14; -11; -10; -9; -7; -6; -5\}$
\task $\{-6; -3; -2; -1; 1; 2; 3; 6\}$
\end{tasks}
\dapan{A}
\end{cauhoi}
\begin{cauhoi}{Câu 9}{8H2-C6-B24-D02-TN-VD-009}
Cho $P=\dfrac{x+4}{x+7}$. Tìm các số nguyên $x$ để $P$ có giá trị nguyên.
\begin{tasks}(2)
\task $\{-10; -8; -6\}$
\task $\{-10; -8; -6; -4\}$
\task $\{-8; -6; -4\}$
\task $\{-3; -1; 1; 3\}$
\end{tasks}
\dapan{B}
\end{cauhoi}
\begin{cauhoi}{Câu 10}{8H2-C6-B24-D02-TN-VD-010}
Cho $P=\dfrac{x+2}{x+3}$. Tìm các số nguyên $x$ để $P$ có giá trị nguyên.
\begin{tasks}(2)
\task $\{-4; -2\}$
\task $\{-1; 1\}$
\task $\{-2\}$
\task $\{-4\}$
\end{tasks}
\dapan{A}
\end{cauhoi}
\begin{cauhoi}{Câu 11}{8H2-C6-B24-D02-TN-VD-011}
Cho $P=\dfrac{x+1}{x+3}$. Tìm các số nguyên $x$ để $P$ có giá trị nguyên.
\begin{tasks}(2)
\task $\{-5; -4; -2\}$
\task $\{-2; -1; 1; 2\}$
\task $\{-4; -2; -1\}$
\task $\{-5; -4; -2; -1\}$
\end{tasks}
\dapan{D}
\end{cauhoi}
\begin{cauhoi}{Câu 12}{8H2-C6-B24-D02-TN-VD-012}
Cho $P=\dfrac{x+8}{x+7}$. Tìm các số nguyên $x$ để $P$ có giá trị nguyên.
\begin{tasks}(2)
\task $\{-1; 1\}$
\task $\{-6\}$
\task $\{-8\}$
\task $\{-8; -6\}$
\end{tasks}
\dapan{D}
\end{cauhoi}
\begin{cauhoi}{Câu 13}{8H2-C6-B24-D02-TN-VD-013}
Cho $P=\dfrac{x+8}{x+2}$. Tìm các số nguyên $x$ để $P$ có giá trị nguyên.
\begin{tasks}(2)
\task $\{-8; -5; -4; -3; -1; 0; 1; 4\}$
\task $\{-6; -3; -2; -1; 1; 2; 3; 6\}$
\task $\{-8; -5; -4; -3; -1; 0; 1\}$
\task $\{-5; -4; -3; -1; 0; 1; 4\}$
\end{tasks}
\dapan{A}
\end{cauhoi}
\begin{cauhoi}{Câu 14}{8H2-C6-B24-D02-TN-VD-014}
Cho $P=\dfrac{x+4}{x+6}$. Tìm các số nguyên $x$ để $P$ có giá trị nguyên.
\begin{tasks}(2)
\task $\{-7; -5; -4\}$
\task $\{-8; -7; -5\}$
\task $\{-2; -1; 1; 2\}$
\task $\{-8; -7; -5; -4\}$
\end{tasks}
\dapan{D}
\end{cauhoi}
\begin{cauhoi}{Câu 15}{8H2-C6-B24-D02-TN-VD-015}
Cho $P=\dfrac{x+4}{x+5}$. Tìm các số nguyên $x$ để $P$ có giá trị nguyên.
\begin{tasks}(2)
\task $\{-4\}$
\task $\{-6; -4\}$
\task $\{-6\}$
\task $\{-1; 1\}$
\end{tasks}
\dapan{B}
\end{cauhoi}
\begin{cauhoi}{Câu 16}{8H2-C6-B24-D02-TN-VD-016}
Cho $P=\dfrac{x+3}{x+6}$. Tìm các số nguyên $x$ để $P$ có giá trị nguyên.
\begin{tasks}(2)
\task $\{-9; -7; -5\}$
\task $\{-3; -1; 1; 3\}$
\task $\{-9; -7; -5; -3\}$
\task $\{-7; -5; -3\}$
\end{tasks}
\dapan{C}
\end{cauhoi}
\begin{cauhoi}{Câu 17}{8H2-C6-B24-D02-TN-VD-017}
Cho $P=\dfrac{x+8}{x+3}$. Tìm các số nguyên $x$ để $P$ có giá trị nguyên.
\begin{tasks}(2)
\task $\{-5; -1; 1; 5\}$
\task $\{-8; -4; -2; 2\}$
\task $\{-8; -4; -2\}$
\task $\{-4; -2; 2\}$
\end{tasks}
\dapan{B}
\end{cauhoi}
\begin{cauhoi}{Câu 18}{8H2-C6-B24-D02-TN-VD-018}
Cho $P=\dfrac{x+1}{x+2}$. Tìm các số nguyên $x$ để $P$ có giá trị nguyên.
\begin{tasks}(2)
\task $\{-3; -1\}$
\task $\{-3\}$
\task $\{-1\}$
\task $\{-1; 1\}$
\end{tasks}
\dapan{A}
\end{cauhoi}
\begin{cauhoi}{Câu 19}{8H2-C6-B24-D02-TN-VD-019}
Cho $P=\dfrac{x+7}{x+2}$. Tìm các số nguyên $x$ để $P$ có giá trị nguyên.
\begin{tasks}(2)
\task $\{-7; -3; -1; 3\}$
\task $\{-7; -3; -1\}$
\task $\{-3; -1; 3\}$
\task $\{-5; -1; 1; 5\}$
\end{tasks}
\dapan{A}
\end{cauhoi}
\begin{cauhoi}{Câu 20}{8H2-C6-B24-D02-TN-VD-020}
Cho $P=\dfrac{x+6}{x+8}$. Tìm các số nguyên $x$ để $P$ có giá trị nguyên.
\begin{tasks}(2)
\task $\{-10; -9; -7; -6\}$
\task $\{-9; -7; -6\}$
\task $\{-10; -9; -7\}$
\task $\{-2; -1; 1; 2\}$
\end{tasks}
\dapan{A}
\end{cauhoi}
\begin{cauhoi}{Câu 21}{8H2-C6-B24-D02-TN-VD-021}
Cho $P=\dfrac{x+3}{x+7}$. Tìm các số nguyên $x$ để $P$ có giá trị nguyên.
\begin{tasks}(2)
\task $\{-9; -8; -6; -5; -3\}$
\task $\{-4; -2; -1; 1; 2; 4\}$
\task $\{-11; -9; -8; -6; -5\}$
\task $\{-11; -9; -8; -6; -5; -3\}$
\end{tasks}
\dapan{D}
\end{cauhoi}
\begin{cauhoi}{Câu 22}{8H2-C6-B24-D02-TN-VD-022}
Cho $P=\dfrac{x+7}{x+4}$. Tìm các số nguyên $x$ để $P$ có giá trị nguyên.
\begin{tasks}(2)
\task $\{-3; -1; 1; 3\}$
\task $\{-7; -5; -3\}$
\task $\{-7; -5; -3; -1\}$
\task $\{-5; -3; -1\}$
\end{tasks}
\dapan{C}
\end{cauhoi}
\begin{cauhoi}{Câu 23}{8H2-C6-B24-D02-TN-VD-023}
Cho $P=\dfrac{x+7}{x+8}$. Tìm các số nguyên $x$ để $P$ có giá trị nguyên.
\begin{tasks}(2)
\task $\{-9; -7\}$
\task $\{-1; 1\}$
\task $\{-9\}$
\task $\{-7\}$
\end{tasks}
\dapan{A}
\end{cauhoi}
\begin{cauhoi}{Câu 24}{8H2-C6-B24-D02-TN-VD-024}
Cho $P=\dfrac{x+5}{x+7}$. Tìm các số nguyên $x$ để $P$ có giá trị nguyên.
\begin{tasks}(2)
\task $\{-2; -1; 1; 2\}$
\task $\{-9; -8; -6; -5\}$
\task $\{-9; -8; -6\}$
\task $\{-8; -6; -5\}$
\end{tasks}
\dapan{B}
\end{cauhoi}
\begin{cauhoi}{Câu 25}{8H2-C6-B24-D02-TN-VD-025}
Cho $P=\dfrac{x+4}{x+2}$. Tìm các số nguyên $x$ để $P$ có giá trị nguyên.
\begin{tasks}(2)
\task $\{-2; -1; 1; 2\}$
\task $\{-3; -1; 0\}$
\task $\{-4; -3; -1; 0\}$
\task $\{-4; -3; -1\}$
\end{tasks}
\dapan{C}
\end{cauhoi}
\begin{cauhoi}{Câu 26}{8H2-C6-B24-D02-TN-VDC-026}
Cho $P=\dfrac{x+1}{x+4}$. Tìm các số nguyên $x$ để $P$ có giá trị nguyên.
\begin{tasks}(2)
\task $\{-7; -5; -3; -1\}$
\task $\{-5; -3; -1\}$
\task $\{-7; -5; -3\}$
\task $\{-3; -1; 1; 3\}$
\end{tasks}
\dapan{A}
\end{cauhoi}
\begin{cauhoi}{Câu 27}{8H2-C6-B24-D02-TN-VDC-027}
Cho $P=\dfrac{x+8}{x+6}$. Tìm các số nguyên $x$ để $P$ có giá trị nguyên.
\begin{tasks}(2)
\task $\{-8; -7; -5\}$
\task $\{-2; -1; 1; 2\}$
\task $\{-7; -5; -4\}$
\task $\{-8; -7; -5; -4\}$
\end{tasks}
\dapan{D}
\end{cauhoi}
\begin{cauhoi}{Câu 28}{8H2-C6-B24-D02-TN-VDC-028}
Cho $P=\dfrac{x+1}{x+8}$. Tìm các số nguyên $x$ để $P$ có giá trị nguyên.
\begin{tasks}(2)
\task $\{-15; -9; -7; -1\}$
\task $\{-9; -7; -1\}$
\task $\{-7; -1; 1; 7\}$
\task $\{-15; -9; -7\}$
\end{tasks}
\dapan{A}
\end{cauhoi}
\begin{cauhoi}{Câu 29}{8H2-C6-B24-D02-TN-VDC-029}
Cho $P=\dfrac{x+5}{x+2}$. Tìm các số nguyên $x$ để $P$ có giá trị nguyên.
\begin{tasks}(2)
\task $\{-3; -1; 1\}$
\task $\{-5; -3; -1; 1\}$
\task $\{-3; -1; 1; 3\}$
\task $\{-5; -3; -1\}$
\end{tasks}
\dapan{B}
\end{cauhoi}
\begin{cauhoi}{Câu 30}{8H2-C6-B24-D02-TN-VDC-030}
Cho $P=\dfrac{x+1}{x+7}$. Tìm các số nguyên $x$ để $P$ có giá trị nguyên.
\begin{tasks}(2)
\task $\{-6; -3; -2; -1; 1; 2; 3; 6\}$
\task $\{-13; -10; -9; -8; -6; -5; -4\}$
\task $\{-13; -10; -9; -8; -6; -5; -4; -1\}$
\task $\{-10; -9; -8; -6; -5; -4; -1\}$
\end{tasks}
\dapan{C}
\end{cauhoi}
\begin{cauhoi}{Câu 31}{8H2-C6-B24-D02-TN-VDC-031}
Cho $P=\dfrac{x+7}{x+6}$. Tìm các số nguyên $x$ để $P$ có giá trị nguyên.
\begin{tasks}(2)
\task $\{-7; -5\}$
\task $\{-7\}$
\task $\{-5\}$
\task $\{-1; 1\}$
\end{tasks}
\dapan{A}
\end{cauhoi}
\begin{cauhoi}{Câu 32}{8H2-C6-B24-D02-TN-VDC-032}
Cho $P=\dfrac{x+5}{x+6}$. Tìm các số nguyên $x$ để $P$ có giá trị nguyên.
\begin{tasks}(2)
\task $\{-7\}$
\task $\{-1; 1\}$
\task $\{-7; -5\}$
\task $\{-5\}$
\end{tasks}
\dapan{C}
\end{cauhoi}
\begin{cauhoi}{Câu 33}{8H2-C6-B24-D02-TN-VDC-033}
Cho $P=\dfrac{x+4}{x+3}$. Tìm các số nguyên $x$ để $P$ có giá trị nguyên.
\begin{tasks}(2)
\task $\{-2\}$
\task $\{-4\}$
\task $\{-1; 1\}$
\task $\{-4; -2\}$
\end{tasks}
\dapan{D}
\end{cauhoi}
\begin{cauhoi}{Câu 34}{8H2-C6-B24-D02-TN-VDC-034}
Cho $P=\dfrac{x+2}{x+4}$. Tìm các số nguyên $x$ để $P$ có giá trị nguyên.
\begin{tasks}(2)
\task $\{-6; -5; -3; -2\}$
\task $\{-6; -5; -3\}$
\task $\{-5; -3; -2\}$
\task $\{-2; -1; 1; 2\}$
\end{tasks}
\dapan{A}
\end{cauhoi}
\begin{cauhoi}{Câu 35}{8H2-C6-B24-D02-TN-VDC-035}
Cho $P=\dfrac{x+8}{x+4}$. Tìm các số nguyên $x$ để $P$ có giá trị nguyên.
\begin{tasks}(2)
\task $\{-4; -2; -1; 1; 2; 4\}$
\task $\{-8; -6; -5; -3; -2; 0\}$
\task $\{-6; -5; -3; -2; 0\}$
\task $\{-8; -6; -5; -3; -2\}$
\end{tasks}
\dapan{B}
\end{cauhoi}
\begin{cauhoi}{Câu 36}{8H2-C6-B24-D02-TN-VDC-036}
Cho $P=\dfrac{x+6}{x+1}$. Tìm các số nguyên $x$ để $P$ có giá trị nguyên.
\begin{tasks}(2)
\task $\{-5; -1; 1; 5\}$
\task $\{-2; 0; 4\}$
\task $\{-6; -2; 0\}$
\task $\{-6; -2; 0; 4\}$
\end{tasks}
\dapan{D}
\end{cauhoi}
\begin{cauhoi}{Câu 37}{8H2-C6-B24-D02-TN-VDC-037}
Cho $P=\dfrac{x+1}{x+5}$. Tìm các số nguyên $x$ để $P$ có giá trị nguyên.
\begin{tasks}(2)
\task $\{-9; -7; -6; -4; -3\}$
\task $\{-7; -6; -4; -3; -1\}$
\task $\{-4; -2; -1; 1; 2; 4\}$
\task $\{-9; -7; -6; -4; -3; -1\}$
\end{tasks}
\dapan{D}
\end{cauhoi}
\begin{cauhoi}{Câu 38}{8H2-C6-B24-D02-TN-VDC-038}
Cho $P=\dfrac{x+3}{x+5}$. Tìm các số nguyên $x$ để $P$ có giá trị nguyên.
\begin{tasks}(2)
\task $\{-2; -1; 1; 2\}$
\task $\{-6; -4; -3\}$
\task $\{-7; -6; -4\}$
\task $\{-7; -6; -4; -3\}$
\end{tasks}
\dapan{D}
\end{cauhoi}
\begin{cauhoi}{Câu 39}{8H2-C6-B24-D02-TN-VDC-039}
Cho $P=\dfrac{x+8}{x+5}$. Tìm các số nguyên $x$ để $P$ có giá trị nguyên.
\begin{tasks}(2)
\task $\{-8; -6; -4; -2\}$
\task $\{-3; -1; 1; 3\}$
\task $\{-6; -4; -2\}$
\task $\{-8; -6; -4\}$
\end{tasks}
\dapan{A}
\end{cauhoi}
\begin{cauhoi}{Câu 40}{8H2-C6-B24-D02-TN-VDC-040}
Cho $P=\dfrac{x+4}{x+8}$. Tìm các số nguyên $x$ để $P$ có giá trị nguyên.
\begin{tasks}(2)
\task $\{-12; -10; -9; -7; -6; -4\}$
\task $\{-4; -2; -1; 1; 2; 4\}$
\task $\{-10; -9; -7; -6; -4\}$
\task $\{-12; -10; -9; -7; -6\}$
\end{tasks}
\dapan{A}
\end{cauhoi}
\begin{cauhoi}{Câu 41}{8H2-C6-B24-D02-TN-VDC-041}
Cho $P=\dfrac{x+7}{x+5}$. Tìm các số nguyên $x$ để $P$ có giá trị nguyên.
\begin{tasks}(2)
\task $\{-7; -6; -4\}$
\task $\{-7; -6; -4; -3\}$
\task $\{-6; -4; -3\}$
\task $\{-2; -1; 1; 2\}$
\end{tasks}
\dapan{B}
\end{cauhoi}
\begin{cauhoi}{Câu 42}{8H2-C6-B24-D02-TN-VDC-042}
Cho $P=\dfrac{x+6}{x+3}$. Tìm các số nguyên $x$ để $P$ có giá trị nguyên.
\begin{tasks}(2)
\task $\{-6; -4; -2; 0\}$
\task $\{-4; -2; 0\}$
\task $\{-3; -1; 1; 3\}$
\task $\{-6; -4; -2\}$
\end{tasks}
\dapan{A}
\end{cauhoi}
\begin{cauhoi}{Câu 43}{8H2-C6-B24-D02-TN-VDC-043}
Cho $P=\dfrac{x+2}{x+5}$. Tìm các số nguyên $x$ để $P$ có giá trị nguyên.
\begin{tasks}(2)
\task $\{-3; -1; 1; 3\}$
\task $\{-8; -6; -4\}$
\task $\{-8; -6; -4; -2\}$
\task $\{-6; -4; -2\}$
\end{tasks}
\dapan{C}
\end{cauhoi}
\begin{cauhoi}{Câu 44}{8H2-C6-B24-D02-TN-VDC-044}
Cho $P=\dfrac{x+2}{x+7}$. Tìm các số nguyên $x$ để $P$ có giá trị nguyên.
\begin{tasks}(2)
\task $\{-12; -8; -6\}$
\task $\{-12; -8; -6; -2\}$
\task $\{-5; -1; 1; 5\}$
\task $\{-8; -6; -2\}$
\end{tasks}
\dapan{B}
\end{cauhoi}
\begin{cauhoi}{Câu 45}{8H2-C6-B24-D02-TN-VDC-045}
Cho $P=\dfrac{x+7}{x+1}$. Tìm các số nguyên $x$ để $P$ có giá trị nguyên.
\begin{tasks}(2)
\task $\{-7; -4; -3; -2; 0; 1; 2; 5\}$
\task $\{-7; -4; -3; -2; 0; 1; 2\}$
\task $\{-4; -3; -2; 0; 1; 2; 5\}$
\task $\{-6; -3; -2; -1; 1; 2; 3; 6\}$
\end{tasks}
\dapan{A}
\end{cauhoi}
\begin{cauhoi}{Câu 46}{8H2-C6-B24-D02-TN-VDC-046}
Cho $P=\dfrac{x+1}{x+6}$. Tìm các số nguyên $x$ để $P$ có giá trị nguyên.
\begin{tasks}(2)
\task $\{-11; -7; -5\}$
\task $\{-7; -5; -1\}$
\task $\{-5; -1; 1; 5\}$
\task $\{-11; -7; -5; -1\}$
\end{tasks}
\dapan{D}
\end{cauhoi}
\begin{cauhoi}{Câu 47}{8H2-C6-B24-D02-TN-VDC-047}
Cho $P=\dfrac{x+8}{x+1}$. Tìm các số nguyên $x$ để $P$ có giá trị nguyên.
\begin{tasks}(2)
\task $\{-2; 0; 6\}$
\task $\{-7; -1; 1; 7\}$
\task $\{-8; -2; 0\}$
\task $\{-8; -2; 0; 6\}$
\end{tasks}
\dapan{D}
\end{cauhoi}
\begin{cauhoi}{Câu 48}{8H2-C6-B24-D02-TN-VDC-048}
Cho $P=\dfrac{x+6}{x+4}$. Tìm các số nguyên $x$ để $P$ có giá trị nguyên.
\begin{tasks}(2)
\task $\{-2; -1; 1; 2\}$
\task $\{-6; -5; -3\}$
\task $\{-6; -5; -3; -2\}$
\task $\{-5; -3; -2\}$
\end{tasks}
\dapan{C}
\end{cauhoi}
\begin{cauhoi}{Câu 49}{8H2-C6-B24-D02-TN-VDC-049}
Cho $P=\dfrac{x+5}{x+8}$. Tìm các số nguyên $x$ để $P$ có giá trị nguyên.
\begin{tasks}(2)
\task $\{-11; -9; -7\}$
\task $\{-3; -1; 1; 3\}$
\task $\{-11; -9; -7; -5\}$
\task $\{-9; -7; -5\}$
\end{tasks}
\dapan{C}
\end{cauhoi}
\begin{cauhoi}{Câu 50}{8H2-C6-B24-D02-TN-VDC-050}
Cho $P=\dfrac{x+3}{x+2}$. Tìm các số nguyên $x$ để $P$ có giá trị nguyên.
\begin{tasks}(2)
\task $\{-3\}$
\task $\{-1\}$
\task $\{-1; 1\}$
\task $\{-3; -1\}$
\end{tasks}
\dapan{D}
\end{cauhoi}
\clearpage\section*{Bảng đáp án trắc nghiệm}
\begin{center}\begin{tabular}{|c|c|c|c|c|c|c|c|c|c|}
\hline
1. D & 2. A & 3. B & 4. B & 5. C & 6. D & 7. D & 8. A & 9. B & 10. A\\\hline
11. D & 12. D & 13. A & 14. D & 15. B & 16. C & 17. B & 18. A & 19. A & 20. A\\\hline
21. D & 22. C & 23. A & 24. B & 25. C & 26. A & 27. D & 28. A & 29. B & 30. C\\\hline
31. A & 32. C & 33. D & 34. A & 35. B & 36. D & 37. D & 38. D & 39. A & 40. A\\\hline
41. B & 42. A & 43. C & 44. B & 45. A & 46. D & 47. D & 48. C & 49. C & 50. D\\\hline
\end{tabular}\end{center}
\end{document}
